% 问题三候选正文：6.1--6.3。
% 本文件只使用 Q3-CONDITIONAL-INTERFACE-v1.0.0 的条件结果，未纳入队友的旧版交接结果。
% 当前为独立候选稿，待科学评审后再决定是否接入 paper/main.tex。

\section{问题三的求解：预算约束下的条件资源配置}
\label{sec:q3}

\subsection{问题分析}
\label{sec:q3-analysis}

\subsubsection{问题要求与求解对象}

问题三要求在给定预算和计算上下文的条件下配置模型规模、数据规模及质量资源，
使预测损失尽可能小。前两问提供的的数据样本并不存在统一的id用于映射、拼接的数据表,其中：
$N$ 和 $D$ 来源于 B1 数据，为标量数据，$Q_B$ 是 B6 数据表中表征质量评分，
而 $p$ 和 $Q_A(p)$ 来自 A附件，主要关于配比的面板数据。当前提供的数据条件难以将$N$ 和 $D$、$Q_B$ 、$p$ 和 $Q_A(p)$ 统一联系在一起考虑。Kaplan 等给出了模型规模与训练损失的幂律刻画，Hoffmann 等进一步在给定训练算力下讨论了
模型规模与训练 token 数的联合配置，本文据此将预算、上下文和质量指标作为条件变量组织资源优化分析。因此本研究将问题三考虑成
\emph{条件资源配置}问题。

\footnote{Kaplan et al., \emph{Scaling Laws for Neural Language Models}, arXiv:2001.08361；
Hoffmann et al., \emph{Training Compute-Optimal Large Language Models}, arXiv:2203.15556。}

求解对象分为三层递进分析：第一层为只使用 B1 的 $N$--$D$ 量建立基线模型 M0；第二层为在
B6加上原生质量评分指标上获得 $N$--$D$--$Q_B$ 模型 M1；第三层为以A提供的配比
$p$ 生成情景，通过显式标注的假设接口建立考虑P的改进的M2模型 。具体而言，对每个配比
$p$ 计算 $Q_A(p)$，依据预先设定的接口情景生成对应的 $Q_B$，再将该 $Q_B$ 代入 M1，
求得相应的 $(N^*,D^*)$，从而形成“配比情景--质量指标--资源配置”的条件面板。


\subsubsection{问题三与前两问的衔接}

前两问提供了三类信息：
\begin{enumerate}
  \item 问题二的 B1 数据提供 $N,D$ 的基线模型；
  \item 问题二的 B6 数据提供能同时考虑$N,D$ 、$Q_B$ )三个标量的数据广义标度率的模型；
  \item 问题一的配比指标提供有限个已观测 $p$ 以及 $Q_A(p)$作为 的情景输入建立考虑$N$ 和 $D$、$Q_B$ 、$p$四量广义标度率的公式。
\end{enumerate}

\subsubsection{求解路线}

采用在 B1数据集求解 考虑获得预算变化下的 $N,D$ M0模型，作为基线；再在 B6数据集内求解 M1模型，比较质量资源与规模资源之间的条件替代关系；最后固定每个 M2/P情景给出的假设 $Q_B$，复用 M1 求解的数据组装成对应的面板数据。

\subsection{数据预分析与变量说明}
\label{sec:q3-data}

\subsubsection{变量、量纲与支持域}

 $N$ 表示模型参数规模（十亿参数），$D$ 表示训练数据规模（十亿
tokens），$L_{\mathrm{ctx}}$ 表示上下文长度，$Q_B$ 表示 B6 原生质量分数。
在 M2/P 中，$p=(p_1,\ldots,p_J)$ 为配比向量，满足非负性以及和为一；问题一使用附件A数据集求解得到的质量分与领域配比分别计算获得来源于A的质量评分指标。
\begin{equation}
  Q_A(p)=\sum_{j=1}^{J}p_jq_j,
  \qquad p_j\geq 0,\qquad \sum_{j=1}^{J}p_j=1.
  \label{eq:q3-qa}
\end{equation}
$Q_A$ 的候选数值仍属于 Q1 的 $0$--$100$ 量尺，$Q_B$ 使用 B6 的
$0.1$--$1.0$ 量尺；

M0 参数拟合采用有界非线性最小二乘模型，以观测 Loss 与标度律预测值的残差平方和为目标，资源配置阶段采用带预算约束和变量范围约束的低维非线性优化模型，最终 求解获得的区间为
\begin{equation}
  N\in[0.070542,11.965825],\qquad
  D\in[0.134,299.893],
  \label{eq:q3-m0-domain}
\end{equation}
其中 $N$ 的单位为十亿参数、$D$ 的单位为十亿 tokens。M1 额外满足

\begin{equation}
  Q_B\in[0.1,1.0],\qquad Q_0=0.1.
  \label{eq:q3-q-domain}
\end{equation}

\subsection{目标函数、预算约束与求解方法}
\label{sec:q3-objective}

\subsubsection{M0：$N$--$D$ 基线优化}

M0 采用 B1 数据集，以有界非线性最小二乘估计方法， Loss 与预测 Loss 的残差平方和为目标拟合广义标度律
\begin{equation}
  \widehat L_{\mathrm{M0}}(N,D)
  =E+A N^{-\alpha}+B D^{-\beta}.
  \label{eq:q3-m0}
\end{equation}
计算所得的参数分别为
\begin{equation*}
\begin{aligned}
  E&=1.6897975629, & A&=0.3539803206, & B&=1.2403055835,\\
  \alpha&=0.3399765819, & \beta&=0.2798781285.
\end{aligned}
\end{equation*}
给定预算 $C$ 和上下文长度 $L_{\mathrm{ctx}}$ 时，M0 的目标为
\begin{equation}
  \min_{N,D}\ \widehat L_{\mathrm{M0}}(N,D)
  \quad\text{s.t.}\quad
  C_{\mathrm{base}}(N,D;L_{\mathrm{ctx}})\leq C,
  \quad (N,D)\in\mathcal D_{\mathrm{B1}}.
  \label{eq:q3-m0-opt}
\end{equation}

\subsubsection{\texorpdfstring{M1：$N$--$D$--$Q_B$ 条件优化}{M1：N-D-QB 条件优化}}

M1 考虑了质量分数$Q_B$ 的因素采用B6数据进行验证和计算。
\begin{equation}
  \widehat L_{\mathrm{M1}}(N,D,Q_B)
  =E+A N^{-\alpha}+B D^{-\beta}+G(1-Q_B).
  \label{eq:q3-m1}
\end{equation}
求解获得的参数分别为
\begin{equation*}
\begin{aligned}
  E&=1.4892660413, & A&=0.5401729883, & B&=1.3167948616,\\
  G&=0.3622474350, & \alpha&=0.2790539239, & \beta&=0.2828029135.
\end{aligned}
\end{equation*}

\subsubsection{成本函数与预算约束}

基础规模成本采用上下文修正的乘积形式，该形式对应 Transformer 训练计算量的常用近似：
模型参数量与训练 token 数将会共同决定基础计算量，前向与反向传播形成系数 $6$；
上下文长度增加会提高注意力相关开销，使用线性修正项 $\eta L_{\mathrm{ctx}}$ 表示。
设$\eta$ 为成本参数，分别采用指数、幂函数和对数渐进成本计算探讨对资源配置结果的影响。
\begin{equation}
  C_{\mathrm{base}}
  =10^{18}N_BD_B\left(6+\eta L_{\mathrm{ctx}}\right),
  \qquad \eta=2\times10^{-4},
  \label{eq:q3-base-cost}
\end{equation}
其中 $N_B,D_B$ 为以十亿单位计的资源变量。质量资源的增量成本写成
\begin{equation}
  C_Q=10^9D_B\left[g(Q_B)-g(Q_0)\right].
  \label{eq:q3-quality-cost}
\end{equation}
为探讨结论对成本形状的敏感性，$g$ 分别取指数、幂函数和对数渐进三种方式进行对比；
总预算约束为
\begin{equation}
  C_{\mathrm{base}}+C_Q\leq C,
  \qquad C\in\{10^{19},10^{22},10^{24}\},
  \qquad L_{\mathrm{ctx}}\in
  \{2048,4096,8192,32768,131072\}.
  \label{eq:q3-constraints}
\end{equation}

\subsubsection{M2/P 配比建模}

对于每个$p$ 先计算 $Q_A(p)$，再按预先登记的接口情景构造一个
假设的 $Q_B$，并将其固定后调用 M1 求解 $(N^*,D^*)$。当前接口包含部分映射
原值、部分映射重归一化以及区间低值、中值和高值等五类情景。该步骤回答的是
接口敏感性问题：不同质量标度假设会使预算前沿如何移动。由于
\begin{equation*}
  \text{formal\_M3\_joint\_fit}=\texttt{false},\qquad
  \text{Q1\_to\_B6\_scale\_calibrated}=\texttt{false},
\end{equation*}
M2/P 的结果不进入 M0 或 M1 的参数再拟合，也不构成正式的
$L(N,D,Q,p)$。

\subsubsection{数值求解与检查}

M0 利用平滑乘积预算约束推导边界解，再用带支持域约束的多起点的SLSQP 检验预算约束、变量范围、收敛状态和目标函数值的一致性，以复核解析推导和边界处理的数值实现。
M1 对输入的正值采用对数参数化，从多个初始点执行 SLSQP。多起点设置用于评估局部求解对初始值的敏感性，并比较不同初始点得到的可行解、目标函数值和边界状态，从而证明模型计算是数值稳定的。每个
解并检查了预算可行性、可行域、收敛态等。M2/P 沿用 M1建立的模型，采用相同的估计质量系数。数据特征呈现出连续变量维度低、目标函数光滑、约束
形式明确，因此优化部分主要采用解析解联合多起点局部优化方式；

\subsubsection{结论分析}
拟合指标与模型分析按参数估计、拟合误差评价、分组留出验证展开：
M0 在 B1 按 $N$ 分组的 8 折留出中，RMSE 为 $0.000102$--$0.000220$，$R^2$ 为
$0.9999992$--$0.9999998$；M1 在 B6 按 $(N,D)$ 单元分组的 5 折留出中，RMSE 为
$0.0512$--$0.0669$，$R^2$ 为 $0.9526$--$0.9705$；B7 的 90 行嵌套扩展诊断中，
RMSE 为 $0.0416$，$R^2=0.9833$。以上拟合数据说明，在考虑了 B1、B6/B7 的数据来源、
Loss 定义、变量量纲和观测范围保持一致的条件下建立的M0，M1拟合模型效果较为精准和稳定。

具体来看，M0 的 8 折留出按照 $N$ 分组，检验了模型在不同参数规模组之间的预测一致性。
RMSE 处于 $10^{-4}$ 数量级，且各折 $R^2$ 均接近 1，说明 B1 数据中的模型规模项、
训练数据项与 Loss 变化具有稳定的幂律关系。因此，M0 适合作为 B1 数据条件下的
$N$--$D$ 资源配置基线，并为后续预算约束优化提供损失函数。

M1 的 5 折留出按照 $(N,D)$ 单元分组，在保留规模组合差异的同时考察 $Q_B$ 质量项的
条件解释能力。RMSE 为 $0.0512$--$0.0669$，$R^2$ 保持在 $0.95$ 以上，说明在 B6
原生质量评分量尺下，模型能够同时描述规模变量和质量变量对 Loss 的主要影响。由此，
M1 可以用于计算质量资源与 $N,D$ 规模资源之间的条件替代关系。

B7 的 90 行数据属于 B6 的嵌套扩展。以 B6 拟合参数计算得到的 RMSE 为 $0.0416$、
$R^2=0.9833$，说明 M1 的条件响应在同一半合成数据构造下对新增实验网格保持较好的
一致性。该结果与 B6 的分组留出结果共同用于评估 M1 的条件适用范围。

从来源诊断看，M0 在 B2、B4、B5 上的 RMSE 分别为 $1.2431$、$0.2927$ 和 $0.1976$，
对应 $R^2$ 分别为 $-5.0728$、$0.6045$ 和 $0.7306$。这些结果体现了不同数据来源
在损失尺度、生成机制和观测范围上的差异，因此问题三的资源优化分别使用 B1 条件下的
M0 和 B6/B7 条件下的 M1，并在各自观测范围与预算情景内解释优化结果。

综合上述结果，M0 主要回答预算约束下 $N$ 与 $D$ 的分配关系，M1 进一步回答加入原生
$Q_B$ 后质量资源与规模资源的权衡，M2/P 则将配比情景转化为 M1 的条件输入。拟合指标
负责说明模型在对应数据条件下的解释能力，优化结果负责呈现预算、上下文和质量成本
变化下的资源配置前沿。
